% Options for packages loaded elsewhere
\PassOptionsToPackage{unicode}{hyperref}
\PassOptionsToPackage{hyphens}{url}
\PassOptionsToPackage{dvipsnames,svgnames,x11names}{xcolor}
\documentclass[
  10pt,
  fontset=fandol]{ctexart}
\usepackage{xcolor}
\usepackage[margin=16mm]{geometry}
\usepackage{amsmath,amssymb}
\setcounter{secnumdepth}{-\maxdimen} % remove section numbering
\usepackage{iftex}
\ifPDFTeX
  \usepackage[T1]{fontenc}
  \usepackage[utf8]{inputenc}
  \usepackage{textcomp} % provide euro and other symbols
\else % if luatex or xetex
  \usepackage{unicode-math} % this also loads fontspec
  \defaultfontfeatures{Scale=MatchLowercase}
  \defaultfontfeatures[\rmfamily]{Ligatures=TeX,Scale=1}
\fi
\usepackage{lmodern}
\ifPDFTeX\else
  % xetex/luatex font selection
\fi
% Use upquote if available, for straight quotes in verbatim environments
\IfFileExists{upquote.sty}{\usepackage{upquote}}{}
\IfFileExists{microtype.sty}{% use microtype if available
  \usepackage[]{microtype}
  \UseMicrotypeSet[protrusion]{basicmath} % disable protrusion for tt fonts
}{}
\makeatletter
\@ifundefined{KOMAClassName}{% if non-KOMA class
  \IfFileExists{parskip.sty}{%
    \usepackage{parskip}
  }{% else
    \setlength{\parindent}{0pt}
    \setlength{\parskip}{6pt plus 2pt minus 1pt}}
}{% if KOMA class
  \KOMAoptions{parskip=half}}
\makeatother
\usepackage{longtable,booktabs,array}
\newcounter{none} % for unnumbered tables
\usepackage{calc} % for calculating minipage widths
% Correct order of tables after \paragraph or \subparagraph
\usepackage{etoolbox}
\makeatletter
\patchcmd\longtable{\par}{\if@noskipsec\mbox{}\fi\par}{}{}
\makeatother
% Allow footnotes in longtable head/foot
\IfFileExists{footnotehyper.sty}{\usepackage{footnotehyper}}{\usepackage{footnote}}
\makesavenoteenv{longtable}
\setlength{\emergencystretch}{3em} % prevent overfull lines
\providecommand{\tightlist}{%
  \setlength{\itemsep}{0pt}\setlength{\parskip}{0pt}}
\usepackage{amsmath,amssymb,mathtools,needspace}
\xeCJKDeclareCharClass{CJK}{"2032,"0400->"04FF,"2160->"216B,"2460->"2473,"25A0,"25C7,"25CB,"25CF}
\IfFontExistsTF{Arial Unicode MS}{\xeCJKsetup{AutoFallBack=true}\setCJKfallbackfamilyfont{\CJKrmdefault}{Arial Unicode MS}}{}
\setcounter{secnumdepth}{0}
\setlength{\emergencystretch}{2em}
\allowdisplaybreaks
\usepackage{bookmark}
\IfFileExists{xurl.sty}{\usepackage{xurl}}{} % add URL line breaks if available
\urlstyle{same}
\hypersetup{
  pdftitle={实用的五点公式},
  colorlinks=true,
  linkcolor={Maroon},
  filecolor={Maroon},
  citecolor={Blue},
  urlcolor={Blue},
  pdfcreator={LaTeX via pandoc}}

\title{实用的五点公式}
\author{}
\date{}

\begin{document}
\maketitle

\subsubsection{4.5.3 实用的五点公式}\label{ux5b9eux7528ux7684ux4e94ux70b9ux516cux5f0f}

设已给出五个节点 \(x_i=x_0+ih,\allowbreak \quad i=0,\allowbreak 1,\allowbreak 2,\allowbreak 3,\allowbreak 4\) 上的函数值，重复同样的手续，不难导出下列五点公式：

\[
m_0=\frac1{12h}[-25f(x_0)+48f(x_1)-36f(x_2)+16f(x_3)-3f(x_4)],
\]

\[
m_1=\frac1{12h}[-3f(x_0)-10f(x_1)+18f(x_2)-6f(x_3)+f(x_4)],
\]

\[
m_2=\frac1{12h}[f(x_0)-8f(x_1)+8f(x_3)-f(x_4)],
\]

\[
m_3=\frac1{12h}[-f(x_0)+6f(x_1)-18f(x_2)+10f(x_3)+3f(x_4)],
\]

\[
m_4=\frac1{12h}[3f(x_0)-16f(x_1)+36f(x_2)-48f(x_3)+25f(x_4)].
\]

其中 \(m_i\) 代表一阶导数 \(f'(x_i)\) 的近似值，读者不难导出这些求导公式的余项。

再记 \(M_i\) 表示二阶导数 \(f''(x_i)\) 的近似值，二阶五点公式如下：

（公式续下页。）

\href{../../source-images/pdf-116.jpeg}{原页图像}

\subsubsection{4.5.3 实用的五点公式（续）}\label{ux5b9eux7528ux7684ux4e94ux70b9ux516cux5f0fux7eed}

\[
M_0=\frac1{12h^2}[35f(x_0)-104f(x_1)+114f(x_2)-56f(x_3)+11f(x_4)],
\]

\[
M_1=\frac1{12h^2}[11f(x_0)-20f(x_1)+6f(x_2)+4f(x_3)-f(x_4)],
\]

\[
M_2=\frac1{12h^2}[-f(x_0)+16f(x_1)-30f(x_2)+16f(x_3)-f(x_4)],
\]

\[
M_3=\frac1{12h^2}[-f(x_0)+4f(x_1)+6f(x_2)-20f(x_3)+11f(x_4)],
\]

\[
M_4=\frac1{12h^2}[11f(x_0)-56f(x_1)+114f(x_2)-104f(x_3)+35f(x_4)].
\]

对于给定的一张数据表，用五点公式求节点上的导数值往往可以获得满意的结果。五个相邻节点的选择原则，一般是在所考察的节点的两侧各取两个邻近的节点，如果一侧的节点数不足两个（即一侧只有一个节点或没有节点），则用另一侧的节点补足。

\textbf{例 4.4}　利用 \(f(x)=\sqrt{x}\) 的一张数据表，按五点公式求节点上的导数值 \(m_i\) 与 \(M_i\)，并与导数的准确值比较，如表 4.8 所示。

\textbf{表 4.8}

{\def\LTcaptype{none} % do not increment counter
\begin{longtable}[]{@{}
  >{\raggedleft\arraybackslash}p{(\linewidth - 10\tabcolsep) * \real{0.1667}}
  >{\raggedleft\arraybackslash}p{(\linewidth - 10\tabcolsep) * \real{0.1667}}
  >{\raggedleft\arraybackslash}p{(\linewidth - 10\tabcolsep) * \real{0.1667}}
  >{\raggedleft\arraybackslash}p{(\linewidth - 10\tabcolsep) * \real{0.1667}}
  >{\raggedleft\arraybackslash}p{(\linewidth - 10\tabcolsep) * \real{0.1667}}
  >{\raggedleft\arraybackslash}p{(\linewidth - 10\tabcolsep) * \real{0.1667}}@{}}
\toprule\noalign{}
\begin{minipage}[b]{\linewidth}\raggedleft
\(x_i\)
\end{minipage} & \begin{minipage}[b]{\linewidth}\raggedleft
\(f(x_i)\)
\end{minipage} & \begin{minipage}[b]{\linewidth}\raggedleft
\(m_i\)
\end{minipage} & \begin{minipage}[b]{\linewidth}\raggedleft
\(f'(x_i)\)
\end{minipage} & \begin{minipage}[b]{\linewidth}\raggedleft
\(M_i/\times10^3\)
\end{minipage} & \begin{minipage}[b]{\linewidth}\raggedleft
\(f''(x_i)/\times10^3\)
\end{minipage} \\
\midrule\noalign{}
\endhead
\bottomrule\noalign{}
\endlastfoot
\(100\) & \(10.000\,000\) & \(0.050\,000\) & \(0.050\,000\) & \(-0.247\,58\) & \(-0.250\,00\) \\
\(101\) & \(10.049\,875\) & \(0.049\,751\) & \(0.049\,752\) & \(-0.245\,91\) & \(-0.246\,30\) \\
\(102\) & \(10.099\,504\) & \(0.049\,507\) & \(0.049\,507\) & \(-0.241\,91\) & \(-0.242\,68\) \\
\(103\) & \(10.148\,891\) & \(0.049\,267\) & \(0.049\,266\) & \(-0.239\,58\) & \(-0.239\,16\) \\
\(104\) & \(10.198\,039\) & \(0.049\,029\) & \(0.049\,029\) & \(-0.236\,91\) & \(-0.235\,72\) \\
\(105\) & \(10.246\,950\) & \(0.048\,795\) & \(0.048\,795\) & \(-0.236\,66\) & \(-0.232\,36\) \\
\end{longtable}
}

\begin{center}\rule{0.5\linewidth}{0.5pt}\end{center}

\subsection{相关原页校注（agent 补充）}\label{ux76f8ux5173ux539fux9875ux6821ux6ce8agent-ux8865ux5145}

以下校注来自本节覆盖的原页，可能也涉及同页相邻小节；原文照录与校正说明分开保留。

\subsubsection{原书 PDF 页 116 的校注}\label{ux539fux4e66-pdf-ux9875-116-ux7684ux6821ux6ce8}

\href{../pages/pdf-116.md}{完整原页转录} · \href{../../source-images/pdf-116.jpeg}{原页图像}

\begin{quote}
校注（agent 补充，原表记法说明）：表 4.8 后两列表头按原扫描保留为 \(M_i/\times10^3\) 与 \(f''(x_i)/\times10^3\)，未改写原表头。表内数值对应导数乘 \(10^3\) 后的数值，例如 \(f''(100)=-1/(4\cdot100^{3/2})=-0.00025\)，表中列为 \(-0.25000\)；此说明仅解释该页缩放记法。\href{../assets/ch04b/review-table-4-8.png}{表 4.8 原扫描局部}。
\end{quote}

\end{document}
